% Ad Familiares 16.7 --- headnote
% ad-familiares-16-07, 17 November 50 BC, Corcyra

\textit{Cicero to Tiro, written from Corcyra on the fifteenth
day before the Kalends of December --- 17 November 50 BC ---
ten days after the three-letter cluster of 16.4--16.6 from
Leucas and Actium. The salutation has narrowed: \textit{Tullius
et Cicero s. d. Tironi suo} --- ``Tullius and Cicero send
greeting to their own Tiro.'' The brother Quintus and his son
no longer share the address line; they are at Buthrotum on the
mainland, and Cicero notes it in the opening sentence. This
closes the second wave of the Tiro-illness cluster from the
homeward leg of the Cilician proconsulship.}

\textit{The party has been weatherbound. Seven days at Corcyra,
no letter from Tiro --- but Cicero will not blame him for the
silence, because the very winds that would carry a letter from
Patrae are the winds that would have let Cicero's own ship
leave the harbour. The reasoning is characteristic: the absence
of news is itself evidence that no news could have come. The
prescription is unchanged from the Leucas letters --- recover
first, then sail when health and season allow --- and the
closing returns to the household-wide affection: \textit{nemo
nos amat qui te non diligat}, ``no one loves us who does not
love you.'' The address \textit{Tiro noster} in the parting
formula is reserved for him alone in book 16.}
